\lecture{2}
\title{How to Train a Neural Net}

\begin{document}

\subsection{Gradient-Based Optimization}

Gradient descent isn't the only gradient-based optimization algorithm.
Even if we can compute the gradient $\mathbf{g} = \nabla_{\theta} J(\theta)$
of a neural net, what should we do with it to best minimize the objective?
\[ \theta^{k + 1} \gets \theta^k + \eta \cdot \textsc{Update}(\mathbf{g}) \]
In general, any optimization algorithm must effectively define
the following two things.
\begin{itemize}
    \item A \textbf{surrogate}---a
    local approximation of the loss function.
    \begin{itemize}
        \item \emph{Example.}
        We might take a first-degree Taylor approximation of the loss.
    \end{itemize}
    \item A \textbf{trust region}---a
    region in which the surrogate is trustable.
    \begin{itemize}
        \item \emph{Example.}
        We might take the trust region to be
        a $\lambda$-radius ball.
    \end{itemize}
\end{itemize}
If we use the two examples given above,
we get exactly gradient descent:
\[ \argmin_{\|\Delta \mathbf{w}\|_2 \le \lambda}
\mathbf{g}^{\top}\Delta \mathbf{w}
= -\lambda \cdot \frac{\mathbf{g}}{\|\mathbf{g}\|_2}. \]
But we can also generalize---if
we take the trust region to be a $\lambda$-radius \emph{square},
we get a different optimizer.
\begin{center}
    \frame{\includegraphics[width=0.25\textwidth]{lecture-02/sign-descent}}
\end{center}
This optimizer is called \textbf{sign descent} because
the minimizer looks like this:
\[ \argmin_{\|\Delta \mathbf{w}\|_{\infty} \le \lambda}
\mathbf{g}^{\top}\Delta \mathbf{w}
= -\lambda \cdot \mathrm{sign}(\mathbf{g}). \]
We can also set the trust region using the
\textbf{spectral norm} of the weight update,
which produces \textbf{spectral descent}.
\[ \argmin_{\|\Delta W\|_{\mathrm{op}} \le \lambda}
\langle G, \Delta W \rangle
= -\lambda \cdot UV^{\top}
~~~ \text{ where } ~~~
G = U\Sigma V^{\top}. \]
\begin{remark}
    Lecture did not give any elaboration on
    what spectral descent is.
    In short, the spectral norm of a matrix
    is its largest \emph{singular} value
    (see \href{/18.701/lectures/19/}{Lecture 19} of 18.701),
    so $-\lambda \cdot UV^{\top}$ is just $-G$
    with every singular value flattened to $\lambda$.
\end{remark}

\subsection{Bells and Whistles}

It turns out that sign descent and spectral descent
(with \emph{bells and whistles}) both yield
commonly-used descent methods.
\begin{center}
    \frame{\includegraphics[width=0.5\textwidth]{lecture-02/descents.png}}
\end{center}
\textbf{Stochasticity.}
The gradient is computed by averaging across
only a subset of the data.
\begin{itemize}
    \item Advantages:
    It's much cheaper,
    and the noise of the stochasticity
    serves as an implicit regularizer.
    \item Disadvantages:
    Higher variance due to gradient approximation.
\end{itemize}
\textbf{Momentum.}
Say the gradient steps are biased
to continue in the direction of the previous update.
\[ \theta^{t + 1} \gets \theta^t - \eta \nabla f(\theta^{t})
+ \alpha \cdot m^t
~~~ \text{ where } ~~~
m^t = \theta^t - \theta^{t - 1} \text{ is the previous update.} \]
Expectedly, the constant $\alpha$ needs to be tuned carefully,
or else advantages may become disadvantages.
\begin{center}
    \frame{\includegraphics[width=0.4\textwidth]{lecture-02/momentum.png}}
\end{center}

\subsection{Gradient Difficulties}

Some functions are easier to optimize than others.
\begin{center}
    \frame{\includegraphics[width=0.5\textwidth]{lecture-02/hard-optimize.png}}
\end{center}
Note that even though most of these functions are not differentiable,
they all have gradients in PyTorch.
Even then, though, PyTorch cannot find a global minimum
when there are vanishing gradients, exploding gradients, or local minima.
\begin{center}
    \frame{\includegraphics[width=0.25\textwidth]{lecture-02/exploding-gradient.png}}
\end{center}
One fix is to use \textbf{gradient clipping}:
if gradients exceed a magnitude $M$,
scale them to have magnitude $M$.
\begin{center}
    \frame{\includegraphics[width=0.25\textwidth]{lecture-02/exploding-fix.png}}
\end{center}
Another fix is to choose a direction not by computing the gradient,
but by considering small perturbations of $\theta$ in all directions
and moving in the direction of lower loss.
\begin{center}
    \frame{\includegraphics[width=0.25\textwidth]{lecture-02/vanishing-fix.png}}
\end{center}

\subsection{Backpropagation}

See \href{/6.790/lectures/13}{Lecture 13} of 6.790
for an explanation of backpropagation
specifically when each layer's computation
$f(\mathbf{x}_{\mathrm{in}}, \theta)$
is a linear layer plus an activation.
We give a more general description here.
(Below, $\theta$ and $\mathbf{W}$ are synonymous.)
\begin{center}
    \frame{\includegraphics[width=0.3\textwidth]{lecture-02/backpropagation.png}}
\end{center}
Notably, for linear layers we have
$f(\mathbf{x}_{\mathrm{in}}, \mathbf{W}) = \mathbf{W}\mathbf{x}_{\mathrm{in}}$,
so it follows that $\mathbf{L}^{\mathbf{x}} = \mathbf{W}$
and $\mathbf{L}^{\theta} = \mathbf{x}_{\mathrm{in}}$.
So backprop is easy.

In practice, not all layers are linear, of course.
For example,
\begin{center}
    \frame{\includegraphics[width=0.5\textwidth]{lecture-02/backprop-example.png}}
\end{center}
But the nonlinear layers like \verb|relu| are also easy to deal with---relu's
Jacobian $\mathbf{L}^{\mathbf{x}}$ is diagonal.

And of course, all of this is compatible with
more general features of DAGs, like merges and branches.
\begin{center}
    \frame{\includegraphics[width=0.4\textwidth]{lecture-02/DAG-features.png}}
\end{center}
Note that \emph{parameter sharing} (used in CNNs, e.g.)
is just a variant of branching.

\end{document}